We are working on this paper.  Below is the general review guidelines:

Classify every issue you find into exactly one of these categories, and lead with the category label and note: (a) does fixing it change the math or logic, (b) does fixing it change the narrative or tone, (c) is it visible to a casual reader vs. only a specialist referee.

* Cohesion Optimization / Modular Consolidation: For any issues point out if it occurs because of a lack of missing [section/paragraph/sentance] solving/answering a general thing and the paper is currently compensating by smearing explanations or repetitive definition across various sections.
* For any issue point out if the fix is actually a simplification here and an improvement elsewhere rather than always adding.
* For any fix suggest distilled addition/changes to prevent bloat and improve clarity.

* FATAL-IDEA (Map failure): The core logic or mathematics is wrong in a way that undermines the paper's central claims. Would cause rejection regardless of presentation.
SHARE (gov failure): Should be fixed before widely sharing as an idea.  Most will interpret the spirit of what is written, but the text needs to be more exact to prevent inaccurate readings.
* CRITICAL-PUBLISH (protocol failure): Does not affect the core logic but would likely cause rejection by a referee or editor. Must be fixed before submission.
* IMPROVE (Mar optimization): Makes the paper clearer or stronger but is not blocking. Worth fixing if the work is not burdensome.
* REMOVE-SIMPLIFY (cap/tol optimization): Anything that is over verbose and can be simplified or should be removed.
POLISH (substrate alignment): Stylistic, LaTeX, or presentational. Fix if convenient, ignore if not.
* SCOPE-DEFENSE: (cap failure) The paper should be stand alone.  Mentioning derivations outside the paper such as in companion papers should be avoided so reviewers can concentrate on reviewing this paper.
* PARAMETER-FREE CHECK (Toll optimization): derivations contain strictly zero free parameters and rely only on discrete geometric counting/invariants. Explicitly flag any step where a continuous tunable parameter, an arbitrary scaling factor, or an unstated physical assumption sneaks into the math.
* ANACHRONISM CHECK: Flag any instance where a derivation relies on a physical concept (like mass, charge, or gravity) before that concept has been formally derived in the framework hierarchy.
* JARGON-CHECK (protocol failure): look for "Jargon Walls." This crosses a number of disciplines of which very few will have all of.  Ideally when using complex systems terminology a bridge or saying it in simpler way when possible ideally (not required if it would explode the length of the paper).
* NAME-CHECK: Informational Energetics or IE refers to the general framework around systems that can persist.  $E_8$-Persistence theory is the name of the physics theory that applies IE to derive the E8 lattice substrate, constants etc
* SHARE (gov failure) The abstract should answer 'what', the introduction should answer 'why', the conclusion should say 'what changed', What do they now know that they didn't before and what comes next
* DIMENSIONALITY AUDIT: Any experimental results should have SI units labeled where appropriate
* IMPROVE (Mar optimization): Any sections overly verbose/edited or repeating previously defined stuff and should be distilled down?
* FLOW: The paper should aim to be organized at every level from the order of the sections to the order of the sentences in a Logical Flow / Narrative Linearity / Monotonic Logic when possible

* This is a complex systems paper. The abstract/conclusion should discuss the complex systems problem of persistence and how IE presents a solution for that. IE is a general framework for all systems that persist and should be stand alone, the validation/example is separate. 
* Complex systems readers should see a universal framework for persistence not just an elaborate setup for a physics derivation.  (AKA the IE section at the start should be cleanly separate from the physics work)
* NARRATIVE-ADMITTANCE: Ensure every mathematical derivation is preceded by a Narrative-Admittance pass. Why does the system/substrate NEED this section? What threat is it solving?
* TONE: This could be thought of framing as a "Dual": (like AdS/CFT). Not "The Standard Model is just a capacity network" 
* TONE: Never declare that standard physics is "arbitrary," "wrong," or an "axiom." Instead, we are showing that standard physics perfectly emerges from this deeper geometric foundation.
* TONE: check for Aggressive/Dismissive Language, replacing any implied critique of standard physics with pure emergent/dual framing.

On the paper itself at a meta level:
* I am currently looking to submit this to nlin.AO (Adaptation and Self-Organizing Systems)
* Like Shannon or E.T. Jaynes this paper is to be built on by everyone, generic in its value, not saying how to optimize a specific system, IE is the abstraction layer
* This is not published yet and we can change anything, a name, the layout, the example, remove something that goes better in a different paper, add something that is critical, literally anything to make the paper better.
* We want to reviewing the general outline of the paper as is.  The high level flow.  Names of sections, order of presentation, name of the paper and section, etc.
* For anything you find suggest latex fixes that can be inserted/edited into the paper as well as if multiple places need to be fixed to fully resolve the issue.
* DISSIPATION-SCRUBBING: Scrub for Verbal Dissipation. Remove repetitions, 'hedging' language (e.g., 'it seems like,' 'one could argue'), and 'filler' adjectives. Every word must perform work to keep the paper within its 'Dimensional Budget'.
* An overarching guage of how complete the paper is: Would a complex systems scientist want to read this in its current form? 
* Check that the citations are correct

The overarching story of the paper:
- The Problem: Structural coherence against entropy
- The Architecture: Six pillars
- The Mechanism: Dissipation density & entropic action 
- The Example / Consequence: The Principle of Least action as survivorship